\prefacechapter
\thispagestyle{empty}

数学始于直觉，但只有当它的假设与论证被明确写出时，它才变得可靠。这份手稿正是从理解这种转变的尝试中生长出来的：熟悉的计算如何被重新组织为定义、定理、证明、算法与模型，并最终能够在不熟悉的情境中反复使用。

本书更适合被称为一份\emph{有引导的二次学习参考}。它主要面向已经接触过一些微积分和线性代数的大学生：他们希望把这些课程同抽象代数、拓扑学、复分析、图论与概率论连接起来理解。本书并不试图替代每一门课程的第一本教材；它的作用是提供一张连贯的地图，给出核心结果的较完整证明、典型例题、应用说明与练习，从而帮助读者在计算能力与结构性理解之间来回移动。

开头关于逻辑与集合论的章节是可选基础，而不是所有读者的入门门槛。以微积分为直接目标的读者，可以从“数学分析 I”开始；对计算与数据更感兴趣的读者，可以从线性代数出发，再进入图论和概率论。基础章节始终保留在那里，供希望考察数学语言与构造如何形式化的读者使用，但本书的结构不再要求每一位读者必须先穿过 ZFC，才能进入分析学。

各部分详略不完全相同，是因为它们在整本书中承担的功能不同。分析与线性代数展开较多，因为它们提供了后续许多章节反复使用的技术。后面的章节更接近选择性的导论，但它们并不只是定理目录。现在，每一章都会说明学习目标，发展有代表性的例题，把理论同科学或计算实践联系起来，并在结尾给出练习与若干提示。

凡是论证机制本身对主题至关重要的地方，证明都会被写出。特别是极限论证、紧致性论证、商构造、收敛定理以及概率中的截断方法，都不会仅仅用“标准证明”一笔带过。若某些结果的完整证明需要远超本章范围的工具，则会明确说明。这个区分很重要：简洁的证明不等于缺失的证明；一个高级定理也不应仅仅因为陈述很短，就被伪装成初等结论。

应用并不是为了装饰理论，而是为了澄清结构。线性代数会联系到最小二乘、数据压缩和主成分分析；抽象代数会联系到编码和密码；多元分析会联系到优化与物理场；拓扑学会联系到数据表示的连续性与商构造；复分析会联系到信号处理与势论；图论会联系到路由、网络和基于图的学习；概率论会联系到统计推断、模拟与随机过程。这些例子有意保持适度：它们的目的在于说明数学语言如何进入真实模型，而不是取代专门的应用课程。

本文在表述上尽量区分个人理解与数学惯例。诸如“线性代数的公理”之类的说法会被替换为更精确的陈述：向量空间公理是一种极其有力的组织观点，但计算的、几何的和算子的观点同样正当。读者应当比较不同表述，并追问在给定问题中哪一种观点最有效。

使用本书的有效方式，是带着纸笔阅读。在读完一个定义后，停下来构造例子与非例子；在阅读证明前，先识别假设，并预判每个假设将在哪里被使用；做章末练习时，不要立刻查看提示。对于计算性主题，可以用一个小例子作数值验证。只有当一个论证能够被重新构造，而不只是被识别出来时，数学理解才会真正变得牢固。

这仍然是一份发展中的手稿。它的目的并不是给出通往本科数学的唯一最终路线，而是提供一条精确、可用、并且其选择对读者可见的路线。

\begin{flushright}
李谨硕 \\
上海交通大学 \\
2026.7.3，于上海
\end{flushright}
